% Part: first-order-logic
% Chapter: tableaux
% Section: proof-theoretic-notions

\documentclass[../../../include/open-logic-section]{subfiles}

\begin{document}

\iftag{FOL}
      {\olfileid{fol}{tab}{ptn}}
      {\olfileid{pl}{tab}{ptn}}
      
\olsection{Begrippen over wat bewijsbaar is}

\begin{editorial}
In deze paragraaf leggen we vast wanneer iets met tableaus bewijsbaar
is en wanneer een verzameling zinnen consistent is.
\end{editorial}

\begin{explain}
We hebben al drie belangrijke begrippen over betekenis gedefinieerd:
geldigheid, logisch gevolg en vervulbaarheid. Nu geven we de
bijbehorende \emph{begrippen over wat bewijsbaar is}. Daarbij kijken we
niet of !!{sentence}s waar zijn in !!{structure}s. We kijken of er
bepaalde gesloten !!{tableau}s bestaan. Een belangrijke ontdekking was
dat beide manieren steeds hetzelfde antwoord geven. Dat is precies wat
de \emph{correctheids-} en \emph{volledigheidsstellingen} zeggen.
\end{explain}


\begin{defn}[Stellingen]
Een !!{sentence}~$!A$ noemen we een \emph{stelling} als er een gesloten
!!{tableau} bestaat voor~$\sFmla{\False}{!A}$. We schrijven
$\Proves !A$ als $!A$ een stelling is en $\Proves/ !A$ als dat niet
zo is.
\end{defn}

\begin{defn}[!!^{derivability}]
!!^a{sentence} $!A$ is precies dan \emph{!!{derivable} uit} een
verzameling !!{sentence}s~$\Gamma$, $\Gamma \Proves !A$, als we een
eindige verzameling
$\{!B_1, \dots, !B_n\} \subseteq \Gamma$ kunnen kiezen waarvoor een
gesloten !!{tableau} bestaat voor de verzameling
\[
\{
  \sFmla{\False}{!A},
  \sFmla{\True}{!B_1}, \dots,
  \sFmla{\True}{!B_n}
  \}.
\]
Als $!A$ niet !!{derivable} is uit $\Gamma$, schrijven we
$\Gamma \Proves/ !A$.
\end{defn}

\begin{defn}[Consistentie]
Een verzameling !!{sentence}s~$\Gamma$ noemen we precies dan
\emph{inconsistent} als we er een eindige verzameling
$\{!B_1, \dots, !B_n\} \subseteq \Gamma$ uit kunnen kiezen waarvoor
een gesloten !!{tableau} bestaat voor de verzameling
\[
\{
  \sFmla{\True}{!B_1}, \dots,
  \sFmla{\True}{!B_n}  
  \}.
\]
Als $\Gamma$ niet inconsistent is, noemen we die verzameling
\emph{consistent}.
\end{defn}

\begin{prop}[Reflexiviteit: wat is aangenomen, volgt ook]
\ollabel{prop:reflexivity}
Als $!A \in \Gamma$, dan $\Gamma \Proves !A$.
\end{prop}

\begin{proof}
Als $!A \in \Gamma$, is $\{!A\}$ een eindige deelverzameling
van~$\Gamma$. Het volgende !!{tableau}
\begin{oltableau}
  [\sFmla{\False}{\formula{A}}, just = \TAss
    [\sFmla{\True}{\formula{A}}, just = \TAss,close]
  ]
\end{oltableau}
is gesloten.
\end{proof}
  

\begin{prop}[Monotoniciteit: extra aannamen mogen erbij]
\ollabel{prop:monotonicity}
Als $\Gamma \subseteq \Delta$ en $\Gamma \Proves !A$, dan $\Delta
\Proves !A$.
\end{prop}

\begin{proof}
Elke eindige deelverzameling van~$\Gamma$ is ook een eindige
deelverzameling van~$\Delta$.
\end{proof}

\begin{prop}[Transitiviteit: afleidingen achter elkaar zetten]
\ollabel{prop:transitivity}
Als $\Gamma \Proves !A$ en $\{!A\} \cup
\Delta \Proves !B$, dan $\Gamma \cup \Delta \Proves !B$.
\end{prop}

\begin{proof}
Stel dat $\{!A\} \cup \Delta \Proves !B$. Dan is er een eindige
deelverzameling $\Delta_0 =
\{!C_1, \dots, !C_n\} \subseteq \Delta$ waarvoor
\begin{align*}
\{\sFmla{\False}{!B}, & \sFmla{\True}{!A}, \sFmla{\True}{!C_1},
\dots, \sFmla{\True}{!C_n}\}
\intertext{een gesloten !!{tableau} bestaat. Stel nu dat
  $\Gamma \Proves !A$. Dan zijn er $!D_1$, \dots,
  $!D_m \subseteq \Gamma$ waarvoor}
\{\sFmla{\False}{!A}, & \sFmla{\True}{!D_1},
\dots, \sFmla{\True}{!D_m}\}
\end{align*}
een gesloten !!{tableau} bestaat.

Bekijk nu het !!{tableau} met de aannamen
\[
\sFmla{\False}{!B},
\sFmla{\True}{!C_1}, \dots, \sFmla{\True}{!C_n},
\sFmla{\True}{!D_1}, \dots, \sFmla{\True}{!D_m}.
\]
Pas de regel \Cut{} toe op~$!A$. Daardoor splitst het !!{tableau} in
twee takken. De ene krijgt $\sFmla{\True}{!A}$; de andere krijgt
$\sFmla{\False}{!A}$. Op de eerste tak staan dan alle elementen van
\[
\{\sFmla{\False}{!B}, \sFmla{\True}{!A},
\sFmla{\True}{!C_1}, \dots, \sFmla{\True}{!C_n}\}
\]
We weten dat er voor deze aannamen een gesloten !!{tableau} bestaat.
Dat kunnen we dus aan deze tak vastmaken. Elke tak die via
$\sFmla{\True}{!A}$ loopt, sluit dan. Op de andere tak staan alle
elementen van
\[
\{\sFmla{\False}{!A}, \sFmla{\True}{!D_1}, \dots,
\sFmla{\True}{!D_m}\}
\]
Ook die kant kunnen we aanvullen tot een gesloten !!{tableau}. Zo
krijgen we één gesloten tableau. Daarmee hebben we laten zien dat
$\Gamma \cup \Delta \Proves !B$.
\end{proof}

Hieruit volgen twee nuttige gevallen. Als $\Gamma \Proves !A$ en $!A
\Proves !B$, dan $\Gamma \Proves !B$. En als $!A_1,
\dots, !A_n \Proves !B$ en $\Gamma \Proves !A_i$ voor elke~$i$, dan
$\Gamma \Proves !B$.

\begin{prop}
\ollabel{prop:incons}
$\Gamma$ is precies dan inconsistent als $\Gamma \Proves !A$ geldt
voor elke !!{sentence}~$!A$.
\end{prop}

\begin{proof}
Opgave.
\end{proof}

\tagprob{FOL}
\begin{prob}
Geef een bewijs van \olref[fol][tab][ptn]{prop:incons}
\end{prob}
\tagendprob

\tagprob{notFOL}
\begin{prob}
Geef een bewijs van \olref[pl][tab][ptn]{prop:incons}
\end{prob}
\tagendprob

\begin{prop}[Compactheid: eindige delen zijn genoeg]
  \ollabel{prop:proves-compact}
  \begin{enumerate}
  \item Als $\Gamma \Proves !A$, dan kunnen we een eindige
    deelverzameling $\Gamma_0 \subseteq \Gamma$ kiezen waarvoor
    $\Gamma_0 \Proves !A$.
  \item Als elke eindige deelverzameling van~$\Gamma$ consistent is,
    dan is $\Gamma$ zelf consistent.
  \end{enumerate}
\end{prop}

\begin{proof}
  \begin{enumerate}
    \item Stel dat $\Gamma \Proves !A$. Dan is er een eindige
      deelverzameling $\Gamma_0 = \{!B_1, \dots, !B_n\}$ en een
      gesloten !!{tableau} voor
      \[
      \{\sFmla{\False}{!A}, \sFmla{\True}{!B_1}, \dots, \sFmla{\True}{!B_n}\}
      \]
      Datzelfde !!{tableau} bewijst ook $\Gamma_0 \Proves !A$.
    \item Stel dat $\Gamma$ inconsistent is. Dan is er een eindige
      deelverzameling $\Gamma_0 = \{!B_1, \dots, !B_n\}$ waarvoor een
      gesloten !!{tableau} bestaat voor
      \[
      \{\sFmla{\True}{!B_1}, \dots, \sFmla{\True}{!B_n}\}
      \]
      Dit gesloten !!{tableau} laat zien dat $\Gamma_0$ inconsistent
      is.
  \end{enumerate}
\end{proof}

\end{document}
